다익스트라 알고리즘은 가중치가 음수인 경우의 문제를 해결할 수 없다는 문제가 있다. 이번 장에서는 음의 가중치가 있어도 단일-전체 최단경로 문제를 해결할 수 있는 벨만-포드 알고리즘을 배울 것이다.

\section{문제 정의}
다음과 같은 최단경로 문제를 생각하자.

\begin{example} \label{ex: Bellman-Ford}
    \textbf{가중치가 음수일 수 있는} 그래프에서 가중치 합이 음수인 사이클이 있는지 판단하고, 없다면 \textbf{1번 정점}에서 시작할 때 각 정점까지의 최단거리를 모두 구하는 프로그램을 작성하라.
\end{example}

가중치가 음수이기 때문에 가중치 합이 음수인 사이클이 존재할 수 있다. 만약 그런 사이클에 도달할 수 있다면 그 사이클을 계속 순환하면 거리가 작아지기 때문에 최단거리가 존재하지 않게 된다. 따라서 음의 가중치가 있는 경우의 단일-전체 최단경로 문제는 가중치 합이 음수인 사이클 존재 여부 판별이 우선되어야 한다.
벨만-포드 알고리즘을 이용하면 좀 더 일반적인 단일-전체 최단경로 문제를 해결할 수 있다.

\section{벨만-포드 알고리즘}
다음과 같은 관찰을 기반으로 한다.

\begin{obs}
    가중치 합이 음수인 사이클이 없다면 각 간선이 $N$번 이상으로 사용되는 최단경로는 존재할 수 없다. 만약 그렇다면 경로 상에 같은 정점이 여러 번 등장할 수 밖에 없고, 이는 가중지 합이 음수인 사이클의 존재를 나타내기 때문이다.
\end{obs}

이제 우리는 모든 정점에 대해 인접한 정점까지 거리를 점화식으로 계산하면 된다. 이 과정의 횟수가 $N$번으로 제한되므로 $O(NM)$에 문제를 해결할 수 있다. \Cref{code: Bellman-Ford}는 구현 코드이다.

\begin{code} \label{code: Bellman-Ford}
벨만-포드 알고리즘 구현 코드
\begin{lstlisting}
#include <bits/stdc++.h>
using namespace std;
typedef pair<int, int> pii;

int n, m;
vector<pii> adj[101010];

int dis[101010];

int main() {
    ios_base::sync_with_stdio(false);
    cin.tie(0);
    cin >> n >> m;
    for (int i = 1; i <= m; i++) {
        int u, v, w;
        cin >> u >> v >> w;
        adj[u].push_back({v, w});
        adj[v].push_back({u, w});
    }
    
    for (int i = 1; i <= n; i++) dis[i] = 1e9;
    dis[1] = 0;
    for (int i = 1; i <= n; i++) {
        for (int v = 1; v <= n; v++) {
            for (int j = 0; j < adj[v].size(); j++) {
                int u = adj[v][j].first, w = adj[v][j].second;
                if (i == n && dis[u] > dis[v] + w) {
                    // there is a nagative cycle
                    return 1;
                }
                dis[u] = min(dis[u], dis[v] + w);
            }
        }
    }
    // there is not any nagative cycle
    // dis is answer
}
\end{lstlisting}
\end{code}

\section{연습문제 A}
\begin{problem}
    \Cref{code: Bellman-Ford}에서 경로를 역추적하는 코드를 추가하라.
\end{problem}
\begin{problem}
    가중치 값이 제약이 없는 그래프에서 1번 정점에서 시작할 때 최장 경로를 존재하는지 판별하고 구하는 프로그램을 작성하라.
\end{problem}

\section{연습문제 B}

\subsection{기초문제}
\begin{problem} \url{https://www.acmicpc.net/problem/1738} \end{problem}
\begin{problem} \url{https://www.acmicpc.net/problem/1865} \end{problem}
\begin{problem} \url{https://www.acmicpc.net/problem/6002} \end{problem}
\begin{problem} \url{https://www.acmicpc.net/problem/6817} \end{problem}
\begin{problem} \url{https://www.acmicpc.net/problem/10360} \end{problem}
\begin{problem} \url{https://www.acmicpc.net/problem/11463} \end{problem}
\begin{problem} \url{https://www.acmicpc.net/problem/11657} \end{problem}
\begin{problem} \url{https://www.acmicpc.net/problem/13317} \end{problem}
\begin{problem} \url{https://www.acmicpc.net/problem/26019} \end{problem}
\begin{problem} \url{https://www.acmicpc.net/problem/30221} \end{problem}

\subsection{응용문제}
\begin{problem} \url{https://www.acmicpc.net/problem/1219} \end{problem}
\begin{problem} \url{https://www.acmicpc.net/problem/2081} \end{problem}
\begin{problem} \url{https://www.acmicpc.net/problem/3680} \end{problem}
\begin{problem} \url{https://www.acmicpc.net/problem/3860} \end{problem}
\begin{problem} \url{https://www.acmicpc.net/problem/4324} \end{problem}
\begin{problem} \url{https://www.acmicpc.net/problem/5942} \end{problem}
\begin{problem} \url{https://www.acmicpc.net/problem/6141} \end{problem}
\begin{problem} \url{https://www.acmicpc.net/problem/7040} \end{problem}
\begin{problem} \url{https://www.acmicpc.net/problem/7141} \end{problem}
\begin{problem} \url{https://www.acmicpc.net/problem/7144} \end{problem}
\begin{problem} \url{https://www.acmicpc.net/problem/7145} \end{problem}
\begin{problem} \url{https://www.acmicpc.net/problem/7634} \end{problem}
\begin{problem} \url{https://www.acmicpc.net/problem/13271} \end{problem}
\begin{problem} \url{https://www.acmicpc.net/problem/22522} \end{problem}

\subsection{도전문제}
\begin{problem} \url{https://www.acmicpc.net/problem/4109} \end{problem}
\begin{problem} \url{https://www.acmicpc.net/problem/6297} \end{problem}
\begin{problem} \url{https://www.acmicpc.net/problem/7332} \end{problem}
\begin{problem} \url{https://www.acmicpc.net/problem/19153} \end{problem}
\begin{problem} \url{https://www.acmicpc.net/problem/19330} \end{problem}
\begin{problem} \url{https://www.acmicpc.net/problem/27400} \end{problem}